% =========================================================
% Gaps in the Rational Numbers
% =========================================================

\section{Gaps in the Rational Numbers}
\label{sec:gaps-in-q}

\subsection{Definitions and Theorems}

The rational numbers are dense and Archimedean, yet they are not complete.  This means that there
are real-number locations that can be approached arbitrarily closely by rationals but are not
occupied by any rational number.  The most famous example is $\sqrt{2}$.

% ---------------------------------------------------------
% Theorem: No Rational Number Has Square 2
% Label: thm:sqrt2-irrational
% ---------------------------------------------------------
\begin{tcolorbox}[colback=thmbox, colframe=thmborder, arc=2pt,
  left=6pt, right=6pt, top=4pt, bottom=4pt,
  title={\small\textbf{Theorem (No Rational Number Has Square $2$)}},
  fonttitle=\small\bfseries]
\begin{theorem}[No Rational Number Has Square $2$]
\label{thm:sqrt2-irrational}
There is no rational number whose square is $2$.
\end{theorem}
\end{tcolorbox}

\begin{remark*}[Standard quantified statement]
\[
\neg\exists r\in\mathbb{Q}\;\bigl(r^2=2\bigr).
\]
\end{remark*}

\begin{remark*}[Predicate reading]
\[
\operatorname{HasSquareEqualToTwo}(r)\coloneqq (r^2=2),
\qquad
\forall r\in\mathbb{Q}\;\neg\operatorname{HasSquareEqualToTwo}(r).
\]
\end{remark*}

\begin{remark*}[Negated quantified statement]
\[
\exists r\in\mathbb{Q}\;\bigl(r^2=2\bigr).
\]
\end{remark*}

\begin{remark*}[Negation predicate reading]
\[
\exists r\in\mathbb{Q}\;\operatorname{HasSquareEqualToTwo}(r).
\]
\end{remark*}

\begin{remark*}[Failure modes]
The theorem would fail exactly if some rational number were an exact square root of $2$.  The proof shows that any such rational, written in lowest terms, would force both numerator and denominator to be even, which is impossible.
\end{remark*}

\begin{remark*}[Failure mode decomposition]
\[
\neg\bigl(\forall r\in\mathbb{Q}\;\neg(r^2=2)\bigr)
\Longleftrightarrow
\exists r\in\mathbb{Q}\;(r^2=2).
\]
\end{remark*}

\begin{remark*}[Interpretation]
This is the first visible gap in $\mathbb{Q}$.  Rational numbers can approximate $\sqrt{2}$ arbitrarily well, but no rational number lands exactly on it.  The obstruction is arithmetic: squaring a reduced rational cannot produce $2$ without violating the parity structure of the integers.
\end{remark*}

\begin{remark*}[Dependencies]
\hyperref[rem:reduced-rational]{Remark~\ref*{rem:reduced-rational}} for reduced-form representatives of rational numbers, together with elementary parity facts for integers.
\end{remark*}

\begin{remark*}[Proof idea]
Argue by contradiction.  Assume $\sqrt{2}=p/q$ with $p/q$ in lowest terms.  Then $p^2=2q^2$, so $p$ is even; substituting $p=2k$ shows $q$ is even as well.  That contradicts the reduced-form assumption.
\end{remark*}

\begin{remark*}[Proof]
\hyperref[prf:sqrt2-irrational]{Go to proof.}
\end{remark*}

\begin{remark}[The gap at $\sqrt{2}$]
Consider the sets
\[
A = \{r\in\mathbb{Q}: r^2<2\}
\qquad\text{and}\qquad
B = \{r\in\mathbb{Q}: r^2>2\}.
\]
Every rational number lies in one of these sets.  Numbers in $A$ approximate $\sqrt{2}$ from
below and numbers in $B$ approximate it from above.  But there is no rational number that fills the
gap between them.
\end{remark}

\begin{remark}[Two-sided approximation]
For example,
\[
1.4<1.41<1.414<\sqrt{2}<1.415<1.42<1.5.
\]
Both sides can be approximated by rational numbers as closely as we please, but the number in the
middle is not itself rational.
\end{remark}

\subsection{Consequences}

\begin{remark}
This is the decisive failure of $\mathbb{Q}$ as a number system for analysis.  The rational numbers
support arithmetic, order, density, mediants, and powerful approximation theorems.  Even so, they
still miss certain limits.  The real numbers are introduced precisely to repair this defect.
\end{remark}
